\documentclass[a4paper,12pt]{article}

\usepackage[T2A]{fontenc}
\usepackage[utf8]{inputenc}
\usepackage[russian]{babel}

\usepackage{amsmath, amssymb, amsthm, mathtools}
\usepackage{helvet}
\renewcommand{\familydefault}{\sfdefault}
\usepackage{icomma}
\usepackage{geometry}
\usepackage{microtype}
\usepackage{tikz}
\usepackage{tkz-euclide}
\usepackage{pgfplots}
\pgfplotsset{compat=1.18}
\usepackage{tabularx, booktabs, longtable}
\usepackage{enumitem}
\usepackage{hyperref}
\usepackage{parskip}
\usepackage{fancyhdr}
\usepackage{setspace}
\usepackage{xcolor}

\definecolor{accent}{HTML}{008080}
\definecolor{darkaccent}{HTML}{004D4D}
\definecolor{textgray}{HTML}{333333}

\geometry{left=2.5cm, right=2.5cm, top=2.5cm, bottom=2.5cm}
\onehalfspacing

\begin{document}

% Титульный лист
\begin{titlepage}
\centering
\vspace*{2cm}
{\Large\bfseries\color{darkaccent} Чек-лист}\par
\vspace{0.5cm}
{{large Тема: Уравнения с модулем и параметром. Окружности}}\par
\vspace{2cm}
{\large Лёвин Артём Александрович}\par
{\large эксклюзивно для Ромы}\par
\vspace{2cm}
{\today}\par
\vfill
\begin{tikzpicture}
\draw[color=accent, line width=1.5pt] 
(0,0) .. controls (2,1.5) and (4,-0.5) .. (6,0.5) .. controls (8,1.5) and (10,-0.5) .. (12,0);
\end{tikzpicture}
\vspace{1cm}
\end{titlepage}

\section*{Чек-лист: Уравнения с модулем и параметром. Окружности}

\subsection*{Основные понятия}
\begin{itemize}[leftmargin=*]
\item \textbf{Модуль}: $|x| = \begin{cases} x, & x \geq 0 \\ -x, & x < 0 \end{cases}$.
\item \textbf{Уравнение окружности}: $(x-a)^2 + (y-b)^2 = R^2$, где $(a;b)$ --- центр, $R$ --- радиус.
\item \textbf{Раскрытие модуля}: рассматривать случаи для положительного и отрицательного подмодульного выражения.
\item \textbf{Условие касания}: расстояние от центра до прямой равно радиусу.
\end{itemize}

\subsection*{Методы решения}
\begin{enumerate}[leftmargin=*]
\item Раскрывать модуль по определению для каждого случая.
\item Строить каждую часть графика в своей области.
\item Проверять «передачу эстафеты» в точках стыка.
\item Для задач с параметром рассматривать все возможные взаимные расположения фигур.
\item Использовать формулу расстояния от точки до прямой: $d = \frac{|Ax_0 + By_0 + C|}{\sqrt{A^2+B^2}}$.
\end{enumerate}

\subsection*{Типичные ошибки}
\begin{itemize}[leftmargin=*]
\item[$\times$] Забывают рассмотреть все случаи раскрытия модуля.
\item[$\times$] Не проверяют принадлежность точек областям определения.
\item[$\times$] Неправильно определяют центр окружности при наличии минуса перед $x$.
\item[$\times$] Не рассматривают все случаи взаимного расположения фигур.
\item[$\times$] Забывают про симметрию фигуры.
\end{itemize}

\newpage
\section*{Тренировочный блок}

\textbf{Уровень 1 (базовый)}

1. Постройте график уравнения: $|x| + |y| = 4$.

2. Найдите центр и радиус окружности: $(x-2)^2 + (y+3)^2 = 16$.

3. Раскройте модуль: $|x-5|$.

\medskip
\textbf{Уровень 2 (средний)}

4. Постройте график: $|x-3| + |y| = 3$.

5. Найдите все значения параметра $a$, при которых окружность $x^2 + y^2 = a^2$ касается прямой $x = 5$.

6. Решите уравнение: $|2x - 4| = 6$.

\medskip
\textbf{Уровень 3 (выше среднего)}

7. Постройте график уравнения: $\frac{|x-3|}{4} + \frac{|y|}{3} = 1$.

8. При каких значениях $a$ система $\begin{cases} (x-2)^2 + y^2 = a^2 \\ |x| + |y| = 4 \end{cases}$ имеет ровно 4 решения?

9. Найдите расстояние от точки $(0;3)$ до прямой $3x + 4y - 24 = 0$.

\medskip
\textbf{Уровень 4 (сложный)}

10. Найдите все значения параметра $a$, при которых уравнение $(|x|-5)^2 + (y-4)^2 = 4$ и уравнение $(x-2)^2 + (y-3)^2 = a^2$ имеют ровно одну общую точку.


\end{document}
